\section{Literature Review}
\label{sec:lit}

This paper draws on four literatures: quantitative overlapping-generations models of demographic aging and pension reform; life-cycle models with endogenous health capital; studies of health financing and social insurance in Latin America; and work on heterogeneity by sex and education in life-cycle fiscal models. The review concentrates on the first, second, and fourth.

%% --------------------------------------------------------------------
\subsection{Demographic Aging and Pension Reform}
\label{subsec:lit_pension}

Since \citet{AuerbachKotlikoff1987}, the multi-period overlapping-generations model with a government sector has been the standard tool for studying the fiscal consequences of demographic change. \citet{KruegerLudwig2007} show that population aging raises the capital--labor ratio, so wages rise and the return on capital falls. The welfare consequences then depend on the composition of a household's resources. Households whose lifetime resources consist mainly of labor income gain from the higher wage, while households that already hold assets lose from the lower return. \citet{KruegerLudwig2007} also show that the adjustment of pension contributions can offset part of the direct effect of aging. \citet{CarvalhoFerreroNechio2016} separate the longevity channel from the dependency-ratio channel in the determination of the equilibrium return; the experiments in this paper change population growth only and hold survival schedules fixed. \citet{AuclertEtAl2021} quantify these forces across countries and find that the sign and size of the aging effect depend on country characteristics. This paper studies the same channels in a Mexican calibration in which the pay-as-you-go contribution rate adjusts endogenously to balance the pension budget.

\citet{Kitao2014} compares four adjustments that restore the balance of a pay-as-you-go system under aging: a higher payroll tax, a lower replacement rate, a higher retirement age, and means-tested benefits. This paper considers the first three. The flat benefit of Section~\ref{sec:model} cannot represent means-testing, so that option is omitted because of the benefit rule rather than a judgment about the instrument. \citet{HuggettVentura1999} compare pension arrangements across stationary economies and show that the choice among them is largely distributional; their exercise is the precedent for the comparison of steady states in this paper. \citet{ConesaKrueger1999} and \citet{NishiyamaSmetters2007} show that an unfunded pillar insures households against uninsurable idiosyncratic wage risk, and \citet{NishiyamaSmetters2007} find that the efficiency gains from reducing it shrink or reverse once transitional generations are compensated. Newborn welfare gains measured across steady states therefore provide an upper bound on the value of a reform once the transition is included.

Differences in mortality by socioeconomic status also affect the incidence of pension systems. \citet{SanchezRomeroEtAl2020} show that longevity gaps by socioeconomic status transfer resources toward the long-lived and reduce the actuarial fairness of pension systems. \citet{Bagchi2019} finds that optimal pension progressivity is insensitive to wealth-dependent mortality. This paper reports the incidence of each adjustment by education group rather than an optimal policy.

The welfare comparisons rest on the condition stated by \citet{Aaron1966}. Where the market return exceeds the growth rate of the contribution base, mandatory participation in an unfunded pillar transfers resources away from participants. \citet{AbelEtAl1989} show that this comparison is more delicate under uncertainty. In the deterministic aggregate environment of this model, the annual return ranges from $5.89$ to $8.58$ percent against base growth of at most $3.51$ percent, so the condition holds in every economy considered. The factor-price channel is a closed-economy object. \citet{AttanasioKitaoViolante2007} show that capital mobility between regions aging at different rates mitigates the factor-price effects of aging, so the closed-economy assumption matters for the factor-price results reported below. \citet{PessinoTerMinassian2021} project the pension and health spending pressure that aging implies for Latin America over the coming decades.

%% --------------------------------------------------------------------
\subsection{Health Capital and the Life Cycle}
\label{subsec:lit_longevity}

The health block follows \citet{Grossman1972}. Health is a depreciable capital stock, medical spending is investment in that stock, and the stock affects survival and, in this paper, labor productivity. Several models build on this structure. \citet{DalgaardStrulik2014} derive aging and death from the interaction of health investment with biological depreciation. \citet{KuhnFeichtinger2015} analyze health and retirement decisions jointly when health capital depreciates with age and can be partially restored. \citet{Fioroni2010} shows that the interaction of precautionary saving with health investment produces consumption and saving profiles that differ from those of models with exogenous health. \citet{HallJones2007} show that when the marginal value of additional life rises with consumption, the income elasticity of health spending exceeds one and the expenditure share of health grows with development, and \citet{MurphyTopel2006} develop the same valuation structure and its empirical magnitudes. The multiplicative health amenity in this paper, and the income gradient in the willingness to pay for health that it generates, follow these two papers.

\citet{JungTran2016} and \citet{JungTranChambers2017} build stochastic general-equilibrium overlapping-generations economies with endogenous health capital and health insurance, and \citet{JungTranChambers2017} use that model for aging experiments for the United States. Their model is the closest to the one in this paper. \citet{Zhao2014} shows that social security itself alters health investment and longevity. \citet{FrankovicWrzaczek2020} trace medical progress through capital accumulation and factor prices, \citet{YewZhang2018} introduce longevity externalities, \citet{Lim2020} shows that public health insurance changes aggregate saving in Korea, and \citet{Carrasco2026} models health investment as a channel through which longer lives raise life-cycle saving and move the equilibrium interest rate. A separate strand treats health as a source of expenditure risk. \citet{DeNardi2010} show that uncertain medical expenses are a primary driver of late-life saving among older American households. \citet{HosseiniKopeckyZhao} estimate a health-deficit process on United States panel data in which health is observed from early adulthood. \citet{GertlerGruber2002} show for a developing economy that major illness is a large and poorly insured shock to household consumption. In this model, the shocks hit the health stock rather than the household budget.

These models are calibrated or estimated mostly on high-income data, in economies where coverage is broad and out-of-pocket payments are a small share of health spending. Such data do not separate spending that builds the health stock from spending that restores it after a shock. \citet{Ozkan2025} separates preventive from curative spending in a life-cycle health-capital model on United States data and documents how their income gradients differ over the life cycle. \citet{GalamaVanKippersluis2013} discuss reformulations of the Grossman model that address the standard criticisms of the framework. For Mexico, \citet{CortesEtAl2024} develop a partial-equilibrium overlapping-generations model in which agents differ by education and face stochastic health transitions, health-dependent productivity, and age-varying out-of-pocket expenditure. The present paper combines endogenous health capital, heterogeneity by sex and education, and an endogenously balanced pay-as-you-go system in a general-equilibrium model calibrated to Mexico and disciplined against its own ENIGH and ENASEM moments. Medical spending in this model is chosen before the health shock is realized, so it is preventive in timing, while most out-of-pocket medical spending in Mexico is curative. Section~\ref{subsec:health_block_evidence} and Appendix~\ref{subsec:health_evidence} examine the consequences of this difference.

%% --------------------------------------------------------------------
\subsection{Health Financing and Social Insurance in Latin America}
\label{subsec:lit_health}

Health shocks are large and poorly insurable through private markets, so pooling arrangements can reduce the associated welfare losses \citep{Arrow1963}. Most Latin American countries operate hybrids of Bismarckian systems, which link coverage to formal employment, and Beveridgean systems, which link it to residency, often with fragmentation and unequal benefit packages \citep{Levy2008}. \citet{Summers1989} shows that mandated contributions act as a tax on formal labor when workers do not value the benefits they finance. \citet{WorldBankInformalizing2025} measure this logic as a formalization tax rate, the net fiscal cost of formal employment inclusive of forgone non-contributory benefits, and find that it is highest for low-income workers. \citet{BaezaPackard2006} document that workers value health coverage, which protects against immediate expenditure risk, more than pension benefits that arrive far in the future. In this model, pension benefits are a flat transfer indexed to average earnings and financed by payroll contributions, with no link between individual contributions and benefits and no formal/informal margin. The results therefore concern pension financing under a stylized pay-as-you-go system rather than formalization incentives, although the contribution rate here is determined by the system's budget rather than imposed.

%% --------------------------------------------------------------------
\subsection{Heterogeneity by Sex and Education}
\label{subsec:lit_gender}

Estimated life-cycle models increasingly study how social security affects men and women differently. \citet{BorellaDeNardiYang2023} estimate a model of couples and singles with survival risk, medical expenses, and the spousal provisions of United States Social Security, and show that those provisions have large and unequal effects on women's labor supply and welfare. \citet{FehrKindermann2018} model a two-person household in which female labor supply responds to childcare, parental leave, and joint versus individual taxation, in an environment without health or endogenous survival. Combining such a household with continuous health capital, health-dependent survival, and endogenous factor prices would raise the computational burden considerably. This paper instead uses four types, two education groups for each sex, each with its own measured age-efficiency profile and baseline survival schedule. These schedules capture in reduced form the differences in labor-market histories and mortality that a structural household model would generate endogenously. Welfare results are reported for all four groups in Section~\ref{sec:experiments}. A structural household model remains a natural extension.

\bigskip
Section~\ref{sec:model} presents the model, which extends \citet{CortesEtAl2024} to general-equilibrium factor prices, a continuous health stock with age-dependent depreciation, and an endogenously balanced pay-as-you-go system. Section~\ref{sec:calibration} describes the calibration and Section~\ref{sec:experiments} the experiments.
